\Article{5}
For brevity we set
\[
\begin{aligned}
(aa+bb+cc)\mathrm{d}t^2+2(aa'+bb'+cc')\mathrm{d}t\mathord{.}\mathrm{d}u&\\
+(a'a'+b'b'+c'c')\mathrm{d}u^2&=\omega
\end{aligned}
\]
\noindent and observe that the differential equation \(\omega=0\) will allow \textit{two} integrations. For when the trinomial \(\omega\) is split into two factors linear in
\SourceLeaf{201}
\(\mathrm{d}t\) and \(\mathrm{d}u\), either one factor or the other must become \(=0\), which will give two different integrations. One integration will correspond to the equation
\[
\begin{aligned}
0={}&(aa+bb+cc)\mathrm{d}t\\
&+\bigl\{aa'+bb'+cc'+i\BareRoot{\bigl((aa+bb+cc)(a'a'+b'b'+c'c')}\\
&\qquad-(aa'+bb'+cc')^2\bigr)\bigr\}\mathrm{d}u
\end{aligned}
\]
\noindent (where \(i\) is written for brevity in place of \(\BareRoot{-1}\), since it is easy to see that the irrational part of the expression must become imaginary); the other will correspond to an entirely similar equation with \(i\) replaced by \(-i\). If, then, the integral of the former equation is
\[
p+iq=\text{Const.}
\]
\noindent where \(p\) and \(q\) denote real functions of \(t\) and \(u\), the other integral will be
\[
p-iq=\text{Const.}
\]
\noindent and the nature of the problem requires that
\[
(\mathrm{d}p+i\,\mathrm{d}q)\mathord{.}(\mathrm{d}p-i\,\mathrm{d}q)\quad\text{or}\quad\mathrm{d}p^2+\mathrm{d}q^2
\]
\noindent be a factor of \(\omega\), or that
\[
\omega=n(\mathrm{d}p^2+\mathrm{d}q^2)
\]
\noindent hold, where \(n\) will be a finite function of \(t\) and \(u\).

Let us consider the trinomial obtained from
\[
\mathrm{d}X^2+\mathrm{d}Y^2+\mathrm{d}Z^2
\]
\noindent by substituting for \(\mathrm{d}X,\mathrm{d}Y,\mathrm{d}Z\) their values in terms of \(T,U,\mathrm{d}T,\mathrm{d}U\). We shall denote this trinomial by \(\Omega\) and assume that, as before, the two integrals of the equation \(\Omega=0\) are
\[
\begin{aligned}
P+iQ&=\text{Const.}\\
P-iQ&=\text{Const.}
\end{aligned}
\]
\noindent and
\[
\Omega=N(\mathrm{d}P^2+\mathrm{d}Q^2)
\]
\noindent where \(P,Q,N\) will denote real functions of \(T\) and \(U\).
\par\indent\SourceLeaf{202}%
Clearly, these integrations can be carried out before solving our main problem (setting aside the general difficulties of integration).

If we now substitute for \(T,U\) functions of \(t,u\) that satisfy the condition of our main problem, \(\Omega\) changes into \(mm\omega\), and we obtain
\[
\frac{(\mathrm{d}P+i\,\mathrm{d}Q)\mathord{.}(\mathrm{d}P-i\,\mathrm{d}Q)}{(\mathrm{d}p+i\,\mathrm{d}q)\mathord{.}(\mathrm{d}p-i\,\mathrm{d}q)}=\frac{mmn}{N}
\]
\noindent It is easy to see, however, that the numerator on the first side of this equation can be divisible by the denominator only if
\[
\begin{aligned}
&\text{either }\mathrm{d}P+i\,\mathrm{d}Q\text{ by }\mathrm{d}p+i\,\mathrm{d}q,\text{ and }\mathrm{d}P-i\,\mathrm{d}Q\text{ by }\mathrm{d}p-i\,\mathrm{d}q,\\
&\text{or }\mathrm{d}P+i\,\mathrm{d}Q\text{ by }\mathrm{d}p-i\,\mathrm{d}q,\text{ and }\mathrm{d}P-i\,\mathrm{d}Q\text{ by }\mathrm{d}p+i\,\mathrm{d}q
\end{aligned}
\]
\noindent the indicated divisibilities hold. In the former case, therefore, \(\mathrm{d}P+i\,\mathrm{d}Q\) will vanish if \(\mathrm{d}p+i\,\mathrm{d}q=0\), or \(P+iQ\) will become constant if \(p+iq\) is taken to be constant; that is, \(P+iQ\) will be a function only of \(p+iq\), and likewise \(P-iQ\) a function of \(p-iq\). In the other case, \(P+iQ\) will be a function of \(p-iq\), and \(P-iQ\) a function of \(p+iq\). It is easy to see that these conclusions also hold conversely: if for \(P+iQ,P-iQ\) we assume functions of \(p+iq,p-iq\) (either respectively or in the reverse order). the finite divisibility of \(\Omega\) by \(\omega\), and hence the proportionality found necessary above, will hold.

It is also easy to see that if, for example,
\[
\begin{aligned}
P+iQ&=f(p+iq)\\
P-iQ&=f'(p-iq)
\end{aligned}
\]
\noindent are set, the nature of the function \(f'\) is already determined by that of \(f\). If the constant quantities that the latter may involve are all real, the other function \(f'\) must be identical with \(f\), so that real values of \(p,q\) always correspond to real values of \(P,Q\); otherwise \(f'\) will differ from \(f\) only in that, in the imaginary elements of \(f\), \(i\) is everywhere replaced by its opposite \(-i\).

We next have
\SourceLeaf{203}
\[
\begin{aligned}
P&=\tfrac{1}{2}f(p+iq)+\tfrac{1}{2}f'(p-iq)\\
iQ&=\tfrac{1}{2}f(p+iq)-\tfrac{1}{2}f'(p-iq)
\end{aligned}
\]
\noindent or, equivalently, with the function \(f\) taken quite arbitrarily (including constant imaginary elements if desired), \(P\) is set equal to the real part and \(iQ\) (in the second solution \(-iQ\)) to the imaginary part of \(f(p+iq)\). From this, by elimination, \(T\) and \(U\) are then expressed as functions of \(t\) and \(u\). This solves the proposed problem quite generally and completely.

